\section{Agent \& Coding: modelの外側にある3つの層}
\sectionkicker{Claude Tag / Grok Build / Agent Quality Flywheel}

\textbf{W32のagent話題を「どのcoding modelが強いか」だけで整理すると、system側の変化を見落とす。}
今号で採用した3候補はいずれもW32以前のartifactだが、並べるとagent systemを支える別々の層が見える。Claude Tagはteam workflowへ常駐する\textbf{product layer}、Grok Buildはtool dispatchやextensionを担う\textbf{orchestration harness}、GoogleのAgent Quality Flywheelは変更を継続的に検証する\textbf{evaluation loop}である。

\begin{tcolorbox}[enhanced,breakable,colback=SurveySoft,colframe=SurveyRule,boxrule=0.4pt,arc=1mm,left=2mm,right=2mm,top=1.5mm,bottom=1.5mm]
\textbf{Agent stackを3層に分けて見る}\par\smallskip
\textbf{Product}: 誰が、どのchannelから、どのdata / toolへagentを接続するか。\par
\textbf{Harness}: context assembly、tool dispatch、subagent、MCP、hookをどう制御するか。\par
\textbf{Evaluation}: 変更後に品質が本当に上がったかを、baselineと再現可能に比較できるか。
\end{tcolorbox}

\subsection*{Claude Tag --- agentがteam contextへ常駐する}

Anthropicは6月23日にClaude Tagを発表し、ClaudeをSlackの選択channelへteam memberとして参加させ、指定したtool・data・codebaseへ接続できる構成を示した。channel内では\texttt{@Claude}へtaskを委任でき、Claudeは関連channel情報からcontextを構築し、将来taskを計画できるとしている\cite{claude-tag-announcement}。

W32での技術的な面白さは新規launchではなく、そのproduct modelが\textbf{governanceの問い}を呼び込んだことにある。Reaction PassではOpen Tagのようなopen / self-host alternativeを介して、model choice、data residency、vendor lock-in、IP / governanceが議論された\cite{claude-tag-x-opentag,claude-tag-x-governance}。persistent team agentでは「どのmodelが答えるか」だけでなく、何を記憶し、どのdataへ常時接続され、誰の権限でtaskを進めるかが運用設計になる。

\subsection*{Grok Build --- modelではなくagent loopそのものを公開する}

SpaceXAIは7月15日、Grok Buildのcoding-agent / TUI sourceをopen source化した。公開範囲にはcontext assembly、model response parsing、tool dispatchからなるagent loop、code read / edit / searchとcommand execution、plan review・inline diffを含むterminal UI、さらにskills、plugins、hooks、MCP servers、subagentsのextension loadingが含まれる\cite{grok-build-open-source}。

重要なのはcoding modelのweightsではなく、\textbf{modelをsoftware engineerとして動かすcontrol planeがinspectableになった}ことである。公式説明ではlocal compileとlocal inferenceへの接続も可能で、harnessを特定model providerと切り離して読む余地が生まれる。一方、W32で観測されたdesktop GUI等のdownstream extensionは今回primary verificationしていないため、本稿ではopen-source harness本体とcommunity派生物を分離する。

\subsection*{Google --- 「良くなった気がする」をevaluation loopへ変える}

Google Developersが6月30日に公開したAgent Quality Flywheelは、agent改善を5段階へ分解する。evaluation dataを準備し、inference / tracesを集め、AutoRaterまたはcustom metricでgradeし、failureを分析し、修正後にbaselineと再比較する\cite{google-agent-quality-flywheel}。

設計上とくに重要なのは、\textbf{optimizerが自分の変更を自分で採点しない}ことだ。Googleはoptimizerとevaluatorを分離し、model-based AutoRaterのabsolute scoreをground truthとは扱わず、run間のdeltaを重視するよう説明している。またsynthetic scenarioはcold-start用であり、成熟後はproduction traceを同じloopへ戻す構成を想定する\cite{google-agent-quality-flywheel}。

\begin{claimboundary}[この3件はW32 launchではない]
Claude Tagは6月23日、Google Flywheel blogは6月30日、Grok Build open sourceは7月15日のeventである。ここで3件を扱う理由はW32の新製品一覧を作るためではなく、W32のagent / coding候補を横断するとproduct・harness・evaluationという異なるsystem layerが同時に重要になっているためである。
\end{claimboundary}

Agent開発の競争軸は、model scoreから徐々にsystem engineeringへ広がっている。teamへ常駐させるなら権限とgovernance、toolを使わせるならharnessの可視性と制御、改善を続けるならevaluationの独立性とregression detectionが必要になる。\textbf{agentの品質は「賢いmodelを選ぶ」だけでは決まらず、その周囲のcontrol loopをどれだけ設計できるかで決まる。}
